\subsection{Quadratic-time coordinate evaluation}

For two distinct vertices $u,w$, define $A(u,w)$ as the balanced average distance between the two components oriented away from the path joining $u$ and $w$.
At an internal vertex, the two outward branches each receive weight $1/2$.
At a leaf, the corresponding distribution is concentrated on that leaf.
Consequently, when the component at $u$ has children $a,b$,
\begin{equation}
 A(u,w)=\tfrac12\{A(a,w)+A(b,w)\}.
 \label{eq:bottom_up}
\end{equation}

To count the work, first fix a target leaf $x$.
The input distances initialize its column: $A(y,x)=D_{yx}$ for every other leaf $y$.
Orient the tree toward $x$ and visit children before their parent.
Each internal vertex has two neighbours farther from $x$, so Equation~\eqref{eq:bottom_up} computes its entry from two values already available.
Each update uses one addition and one multiplication by $1/2$, regardless of how many leaves lie below that vertex.
There are $N-2$ internal vertices, giving $O(N)$ work for the entire column $A(\cdot,x)$.
Repeating this for all $N$ target leaves costs $O(N^2)$.

The SPR evaluator also needs columns indexed by internal vertices.
For an internal target $w$, symmetry supplies every leaf entry $A(x,w)=A(w,x)$ from the columns just computed.
Orienting the tree toward $w$ then permits the same averaging recurrence within each of its three incident branches.
There are $N-3$ internal vertices other than $w$, and the diagonal $A(w,w)$ is unused.
Thus each of the remaining $N-2$ columns also costs $O(N)$, including initialization and traversal:
\[
 \underbrace{N\,O(N)}_{\text{leaf-target columns}}
 +\underbrace{(N-2)\,O(N)}_{\text{internal-target columns}}
 =O(N^2).
\]

The implementation uses a fixed computational rooting to realize these dependencies through bottom-up and top-down passes.
When $u$ is an ancestor of $w$, the component at $u$ away from $w$ combines the sibling branch and the parent branch:
\begin{equation}
 A(u,w)=\tfrac12\{A(\operatorname{sib}_u(w),w)
                         +A(\operatorname{parent}(u),w)\}.
 \label{eq:top_down}
\end{equation}
At the degree-three computational root, the analogous average uses the two branches not containing $w$.
Each column requires $O(N)$ visits across the two passes, including on an unbalanced tree.

Once $A$ is available, a second reuse argument gives all coordinate changes.
For a fixed pruned component, one traversal of the standard BME SPR recurrence propagates score differences across all regraft edges in $O(N)$ time~\cite{hordijk2005spr}.
Each propagation step is an adjacent interchange of four directed components.
Its score change is a linear combination of their balanced distances, with coefficient $1/4$, requiring a constant number of table lookups and arithmetic operations.
Successive steps accumulate the change relative to the original attachment.
Phylo2Vec supplies the mapping from each valid coordinate value to one of these regraft edges.
For row $i$, the traversal therefore supplies all $2i+1$ coordinate entries through $O(i+1)$ lookups.
The explicit mapping, including the collapsed-root boundary, is in Appendix~\ref{sec:mapping}.

\begin{proposition}[Cost and scope of a coordinate sweep]
\label{prop:quadratic}
Given an ordered baseline, its leaf permutation, and constant-time access to the entries of a supplied distance matrix, constructing $A$ and computing all $(N-1)^2$ entries in Equation~\eqref{eq:discrete_gradient} requires $O(N^2)$ arithmetic operations.
The dense implementation uses $O(N^2)$ storage.
In exact arithmetic, these entries equal the corresponding full-tree BME differences.
\end{proposition}
\begin{proof}
There are $2N-2$ columns of $A$, each constructed in $O(N)$ work by the recurrences above.
For coordinate $i$, the two half-tree traversals visit $O(N)$ regraft edges, with constant work per edge, and the coordinate mapping takes $O(i+1)$ work.
Thus the total is
\[
 \underbrace{O(N^2)}_{\text{construct }A}
 +\underbrace{\sum_{i=0}^{N-2}\bigl[O(N)+O(i+1)\bigr]}_{\text{traverse regraft edges and extract each row}}
 =O(N^2).
\]
The dense distance and coordinate tables each have $O(N^2)$ entries.
The recurrence preserves the Pauplin score difference at each traversal step, so selecting a mapped edge gives Equation~\eqref{eq:discrete_gradient}.
\end{proof}

The proposition concerns the coordinate neighborhood of the current labeling.
It does not assert that this neighborhood is the entire SPR neighborhood, or that coordinate descent finds a global BME optimum.
There are $\sum_{i=0}^{N-2}(2i+1)=(N-1)^2$ valid entries.
Any algorithm that explicitly writes one scalar per entry requires $\Omega(N^2)$ work in the standard unit-cost model, so the evaluator matches this output-size lower bound.
This observation does not give a lower bound for returning only the best move, and does not improve the asymptotic sweep cost of classical BME SPR evaluation.

